\section{Kết quả thực nghiệm}
\label{sec:results}

\subsection{Tính đầy đủ của hai lần chạy}

Cả hai cấu hình structural SCM và path ablation đều xử lý thành công toàn bộ
\(250\) câu hỏi. Không có prediction bị thiếu, không có ID nằm ngoài
benchmark và không có pipeline error. Prediction coverage và success rate
đều đạt \(100\%\).

Các kết quả trong phần này chỉ sử dụng hai bộ artifact paired:

\begin{verbatim}
pipeline_predictions_structural_scm_v41_step5_v51.json
pipeline_predictions_path_ablation_v41_step5_v51.json
\end{verbatim}

Hai bộ prediction đều được tạo bởi runner
\texttt{2.1-query-aware-answer} và được đánh giá bằng Step 7
\texttt{2.0-current-pipeline-schema}. Các artifact cũ không có hậu tố
\texttt{step5\_v51} không được đưa vào kết quả chính.

\subsection{Kết quả tổng thể}

Bảng~\ref{tab:results} trình bày các metric chính của cấu hình structural
SCM. Do paired path ablation đạt cùng toàn bộ quality metric, kết quả của
hai cấu hình được so sánh riêng tại
Bảng~\ref{tab:paired-ablation}.

\begin{table}[t]
    \centering
    \caption{Kết quả tổng thể của pipeline structural SCM trên \(250\) câu
    hỏi.}
    \label{tab:results}
    \begin{tabular}{lr}
        \hline
        \textbf{Metric} & \textbf{Kết quả} \\
        \hline
        Prediction coverage & \(100{,}00\%\) \\
        Success rate & \(100{,}00\%\) \\
        Rule Recall@5 & \(69{,}93\%\) \\
        Event Recall@5 & \(74{,}66\%\) \\
        Top-1 Exact Path & \(35{,}20\%\) \\
        Oracle Exact Path & \(71{,}20\%\) \\
        Verification Accuracy & \(92{,}00\%\) \\
        Verification Macro-F1 & \(97{,}73\%\) \\
        Verification Balanced Accuracy & \(95{,}65\%\) \\
        Answer Token F1 & \(68{,}69\%\) \\
        Answer ROUGE-L F1 & \(67{,}67\%\) \\
        Citation F1 & \(62{,}86\%\) \\
        \hline
    \end{tabular}
\end{table}

Kết quả cho thấy retrieval đạt recall tương đối cao ở mức rule và event.
Verification có độ chính xác cao hơn retrieval path và answer generation.
Khoảng cách lớn giữa Top-1 Exact Path và Oracle Exact Path cho thấy path
ranking là một nút thắt quan trọng của pipeline.

\subsection{Kết quả retrieval}

Bảng~\ref{tab:retrieval-results} trình bày recall, precision và hit rate tại
các giá trị \(k\) chính.

\begin{table}[t]
    \centering
    \caption{Kết quả retrieval theo \(k\).}
    \label{tab:retrieval-results}
    \begin{tabular}{lrrr}
        \hline
        \textbf{Metric} & \(\mathbf{k=1}\) & \(\mathbf{k=3}\) &
        \(\mathbf{k=5}\) \\
        \hline
        Rule Recall@\(k\) &
        \(33{,}12\%\) & \(65{,}11\%\) & \(69{,}93\%\) \\
        Rule Precision@\(k\) &
        \(68{,}40\%\) & \(44{,}93\%\) & \(29{,}04\%\) \\
        Rule Hit@\(k\) &
        \(68{,}40\%\) & \(85{,}60\%\) & \(88{,}40\%\) \\
        \hline
        Event Recall@\(k\) &
        \(23{,}37\%\) & \(66{,}83\%\) & \(74{,}66\%\) \\
        Event Precision@\(k\) &
        \(71{,}20\%\) & \(68{,}00\%\) & \(45{,}76\%\) \\
        Event Hit@\(k\) &
        \(71{,}20\%\) & \(88{,}00\%\) & \(89{,}20\%\) \\
        \hline
    \end{tabular}
\end{table}

Tại \(k=1\), precision của cả rule và event tương đối cao, nhưng recall còn
thấp vì mỗi câu hỏi có thể yêu cầu nhiều gold item trên một causal path.
Khi tăng \(k\) lên \(5\), recall tăng đáng kể, trong khi precision giảm do
retriever bổ sung nhiều ứng viên không thuộc gold set.

Các retrieval metric bổ sung gồm:

\begin{itemize}
    \item Rule Macro Recall: \(74{,}41\%\);
    \item Rule Macro Precision: \(12{,}73\%\);
    \item Rule Macro F1: \(21{,}63\%\);
    \item Rule MRR: \(77{,}23\%\);
    \item Rule MAP: \(61{,}41\%\);
    \item Event Macro Recall: \(76{,}78\%\);
    \item Event Macro Precision: \(29{,}50\%\);
    \item Event Macro F1: \(42{,}52\%\);
    \item Event MRR: \(79{,}72\%\).
\end{itemize}

Rule precision tổng thể thấp hơn event precision vì prediction lưu không chỉ
các rule trên primary path mà còn các direct rule hit và rule ứng viên bổ
sung. Kết quả này cho thấy retriever ưu tiên coverage hơn việc tạo một tập
rule tối thiểu.

\subsection{Kết quả causal path}

Bảng~\ref{tab:path-results} trình bày kết quả đánh giá primary path và oracle
path.

\begin{table}[t]
    \centering
    \caption{Kết quả đánh giá causal path.}
    \label{tab:path-results}
    \begin{tabular}{lr}
        \hline
        \textbf{Metric} & \textbf{Kết quả} \\
        \hline
        Top-1 Exact Path Match & \(35{,}20\%\) \\
        Top-1 Hop Accuracy & \(53{,}60\%\) \\
        Top-1 Path Length Accuracy & \(77{,}60\%\) \\
        Top-1 Edge Precision & \(54{,}40\%\) \\
        Top-1 Edge Recall & \(54{,}00\%\) \\
        Top-1 Edge F1 & \(54{,}13\%\) \\
        Top-1 Event F1 & \(61{,}49\%\) \\
        Top-1 Final Event Accuracy & \(49{,}60\%\) \\
        Top-1 Path Rule F1 & \(47{,}14\%\) \\
        \hline
        Oracle Exact Path Match & \(71{,}20\%\) \\
        Oracle Edge F1 & \(75{,}33\%\) \\
        Oracle Event F1 & \(76{,}59\%\) \\
        Oracle Final Event Accuracy & \(73{,}60\%\) \\
        Oracle Path Rule F1 & \(75{,}74\%\) \\
        \hline
        Path Ranking Error Rate & \(36{,}00\%\) \\
        Path Retrieval Error Rate & \(28{,}80\%\) \\
        \hline
    \end{tabular}
\end{table}

Top-1 Exact Path chỉ đạt \(35{,}20\%\), trong khi Oracle Exact Path đạt
\(71{,}20\%\). Chênh lệch tuyệt đối là:

\[
71{,}20\% - 35{,}20\% = 36{,}00
\text{ điểm phần trăm}.
\]

Điều này cho thấy trong \(36\%\) số mẫu, một candidate path khớp gold tốt hơn
đã xuất hiện trong retrieval output nhưng không được chọn làm primary path.
Ngoài ra, \(28{,}80\%\) mẫu không có gold path đầy đủ trong candidate set,
phản ánh lỗi ở giai đoạn retrieval thay vì ranking.

Path Length Accuracy đạt \(77{,}60\%\), cao hơn Exact Path Match. Như vậy,
retriever thường xác định đúng số hop nhưng có thể chọn sai seed event,
mediator, outcome hoặc rule cụ thể.

\subsection{Kết quả xác minh}

Query-aware verifier đạt Accuracy \(92{,}00\%\), Macro-F1 \(97{,}73\%\) và
Balanced Accuracy \(95{,}65\%\). Confusion matrix được trình bày trong
Bảng~\ref{tab:verification-confusion}.

\begin{table}[t]
    \centering
    \caption{Confusion matrix của verification decision.}
    \label{tab:verification-confusion}
    \begin{tabular}{lrrr}
        \hline
        \textbf{Gold} &
        \textbf{SUPPORTED} &
        \textbf{REJECT} &
        \textbf{UNCERTAIN} \\
        \hline
        \texttt{SUPPORTED} & \(210\) & \(0\) & \(20\) \\
        \texttt{REJECT\_DIRECT\_CLAIM} & \(0\) & \(20\) & \(0\) \\
        \texttt{UNCERTAIN} & \(0\) & \(0\) & \(0\) \\
        \hline
    \end{tabular}
\end{table}

Tất cả \(20\) direct-edge negative claims đều được phân loại đúng thành
\texttt{REJECT\_DIRECT\_CLAIM}. Trong \(230\) gold
\texttt{SUPPORTED} sample, hệ thống phân loại đúng \(210\) mẫu và chuyển
\(20\) mẫu thành \texttt{UNCERTAIN}. Không có trường hợp
\texttt{SUPPORTED} bị nhầm thành \texttt{REJECT\_DIRECT\_CLAIM} hoặc ngược
lại.

Recall của hai lớp có gold support là:

\begin{itemize}
    \item \texttt{SUPPORTED}: \(91{,}30\%\);
    \item \texttt{REJECT\_DIRECT\_CLAIM}: \(100{,}00\%\).
\end{itemize}

Precision của cả hai lớp đều đạt \(100\%\). Điều này cho thấy verifier có xu
hướng abstain bằng nhãn \texttt{UNCERTAIN} thay vì đưa ra một quyết định
khẳng định sai. Tuy nhiên, benchmark không có gold sample thuộc lớp
\texttt{UNCERTAIN}, nên chưa thể đánh giá đầy đủ độ chính xác của cơ chế
abstention.

\subsection{Kết quả câu trả lời}

Các metric answer generation được trình bày trong
Bảng~\ref{tab:answer-results}.

\begin{table}[t]
    \centering
    \caption{Kết quả answer generation.}
    \label{tab:answer-results}
    \begin{tabular}{lr}
        \hline
        \textbf{Metric} & \textbf{Kết quả} \\
        \hline
        Token Precision & \(65{,}34\%\) \\
        Token Recall & \(75{,}40\%\) \\
        Token F1 & \(68{,}69\%\) \\
        ROUGE-L F1 & \(67{,}67\%\) \\
        Normalized Exact Match & \(0{,}00\%\) \\
        Answer Present & \(100{,}00\%\) \\
        Trung bình token của prediction & \(48{,}18\) \\
        Trung bình token của gold answer & \(45{,}07\) \\
        \hline
    \end{tabular}
\end{table}

Token Recall cao hơn Token Precision, cho thấy câu trả lời thường bao phủ
phần lớn nội dung gold nhưng chứa thêm causal chain, phần giải thích hoặc
citation. Normalized Exact Match bằng \(0\%\) không đồng nghĩa mọi câu trả
lời đều sai, vì predicted answer thường bổ sung cấu trúc giải thích và marker
\([E_i]\) mà gold answer không chứa.

Token F1 và ROUGE-L F1 có giá trị gần nhau, lần lượt là \(68{,}69\%\) và
\(67{,}67\%\). Điều này cho thấy phần token được bao phủ nhìn chung vẫn giữ
được thứ tự tương đối gần với gold answer.

\subsection{Kết quả citation}

Kết quả citation được trình bày trong Bảng~\ref{tab:citation-results}.

\begin{table}[t]
    \centering
    \caption{Kết quả đánh giá citation.}
    \label{tab:citation-results}
    \begin{tabular}{lr}
        \hline
        \textbf{Metric} & \textbf{Kết quả} \\
        \hline
        Citation Precision & \(62{,}69\%\) \\
        Citation Recall & \(64{,}80\%\) \\
        Citation F1 & \(62{,}86\%\) \\
        Citation Exact Match & \(43{,}20\%\) \\
        \hline
    \end{tabular}
\end{table}

Citation Recall cao hơn Citation Precision một khoảng nhỏ. Hệ thống thường
trích được phần lớn gold article nhưng đôi khi bổ sung article từ một
candidate path khác hoặc bỏ sót một căn cứ nằm trên gold path. Citation
Exact Match đạt \(43{,}20\%\), cho thấy việc khớp hoàn toàn tập căn cứ vẫn là
một bài toán khó hơn so với chỉ tìm được ít nhất một article liên quan.

\subsection{Kết quả theo loại câu hỏi}

Bảng~\ref{tab:question-type-results} trình bày các metric chính theo question
type. Các ký hiệu R@5, E@5, T1, OR, VA, AF1 và CF1 lần lượt là Rule Recall@5,
Event Recall@5, Top-1 Exact Path, Oracle Exact Path, Verification Accuracy,
Answer Token F1 và Citation F1.

\begin{table*}[t]
    \centering
    \scriptsize
    \caption{Kết quả theo loại câu hỏi.}
    \label{tab:question-type-results}
    \begin{tabular}{lrrrrrrrr}
        \hline
        \textbf{Loại câu hỏi} &
        \textbf{N} &
        \textbf{R@5} &
        \textbf{E@5} &
        \textbf{T1} &
        \textbf{OR} &
        \textbf{VA} &
        \textbf{AF1} &
        \textbf{CF1} \\
        \hline
        Branch reasoning &
        \(4\) &
        \(60{,}00\) &
        \(60{,}00\) &
        \(0{,}00\) &
        \(0{,}00\) &
        \(100{,}00\) &
        \(69{,}33\) &
        \(43{,}57\) \\

        Bridge-convergence fallback &
        \(6\) &
        \(100{,}00\) &
        \(100{,}00\) &
        \(83{,}33\) &
        \(100{,}00\) &
        \(100{,}00\) &
        \(65{,}03\) &
        \(91{,}67\) \\

        Bridge event &
        \(40\) &
        \(81{,}25\) &
        \(86{,}67\) &
        \(50{,}00\) &
        \(95{,}00\) &
        \(100{,}00\) &
        \(75{,}95\) &
        \(75{,}42\) \\

        Chain explanation &
        \(35\) &
        \(78{,}57\) &
        \(82{,}86\) &
        \(40{,}00\) &
        \(80{,}00\) &
        \(100{,}00\) &
        \(77{,}22\) &
        \(67{,}62\) \\

        Convergence reasoning &
        \(10\) &
        \(69{,}33\) &
        \(82{,}50\) &
        \(0{,}00\) &
        \(0{,}00\) &
        \(100{,}00\) &
        \(84{,}32\) &
        \(59{,}66\) \\

        Forward multihop &
        \(70\) &
        \(69{,}29\) &
        \(79{,}52\) &
        \(35{,}71\) &
        \(87{,}14\) &
        \(100{,}00\) &
        \(76{,}89\) &
        \(68{,}10\) \\

        Reverse reasoning &
        \(45\) &
        \(46{,}67\) &
        \(42{,}22\) &
        \(8{,}89\) &
        \(26{,}67\) &
        \(60{,}00\) &
        \(32{,}67\) &
        \(36{,}32\) \\

        Yes/no counterexample &
        \(20\) &
        \(82{,}50\) &
        \(88{,}33\) &
        \(55{,}00\) &
        \(95{,}00\) &
        \(100{,}00\) &
        \(89{,}16\) &
        \(67{,}50\) \\

        Yes/no positive &
        \(20\) &
        \(67{,}50\) &
        \(70{,}00\) &
        \(45{,}00\) &
        \(70{,}00\) &
        \(90{,}00\) &
        \(64{,}25\) &
        \(62{,}94\) \\
        \hline
    \end{tabular}
\end{table*}

Các nhóm bridge event, causal chain explanation và forward multihop đạt
Answer Token F1 lần lượt là \(75{,}95\%\), \(77{,}22\%\) và \(76{,}89\%\).
Convergence reasoning đạt Answer Token F1 cao nhất trong các câu hỏi không
phải counterexample, với \(84{,}32\%\).

Yes/no counterexample đạt Answer Token F1 \(89{,}16\%\) và Verification
Accuracy \(100\%\). Kết quả này phản ánh khả năng phân biệt quan hệ trực tiếp
với đường gián tiếp. Tuy nhiên, nhóm này không cung cấp ground truth đầy đủ
cho hard intervention.

Reverse reasoning là nhóm yếu nhất:

\begin{itemize}
    \item Rule Recall@5: \(46{,}67\%\);
    \item Event Recall@5: \(42{,}22\%\);
    \item Top-1 Exact Path: \(8{,}89\%\);
    \item Oracle Exact Path: \(26{,}67\%\);
    \item Verification Accuracy: \(60{,}00\%\);
    \item Answer Token F1: \(32{,}67\%\);
    \item Citation F1: \(36{,}32\%\).
\end{itemize}

Kết quả này cho thấy lỗi reverse reasoning xuất hiện ngay từ retrieval, sau
đó lan truyền sang path selection, verification và answer generation.

Top-1 và Oracle Exact Path đều bằng \(0\%\) trên branch reasoning và
convergence reasoning. Tuy nhiên, các nhóm này vẫn đạt Answer Token F1 lần
lượt là \(69{,}33\%\) và \(84{,}32\%\). Sự khác biệt cho thấy metric path
tuyến tính không phù hợp để đánh giá đầy đủ đầu ra nhiều nhánh hoặc hội tụ.

\subsection{Kết quả theo độ khó}

Bảng~\ref{tab:difficulty-results} trình bày kết quả theo difficulty.

\begin{table}[t]
    \centering
    \caption{Kết quả theo độ khó.}
    \label{tab:difficulty-results}
    \begin{tabular}{lrr}
        \hline
        \textbf{Metric} & \textbf{Medium} & \textbf{Hard} \\
        \hline
        Số câu hỏi & \(181\) & \(69\) \\
        Rule Recall@5 & \(67{,}13\%\) & \(77{,}29\%\) \\
        Event Recall@5 & \(71{,}45\%\) & \(83{,}07\%\) \\
        Top-1 Exact Path & \(34{,}81\%\) & \(36{,}23\%\) \\
        Oracle Exact Path & \(72{,}38\%\) & \(68{,}12\%\) \\
        Verification Accuracy & \(88{,}95\%\) & \(100{,}00\%\) \\
        Answer Token F1 & \(63{,}90\%\) & \(81{,}25\%\) \\
        Citation F1 & \(62{,}02\%\) & \(65{,}04\%\) \\
        \hline
    \end{tabular}
\end{table}

Nhóm hard đạt kết quả answer và verification cao hơn nhóm medium. Đây không
nên được diễn giải rằng độ khó cao làm hệ thống hoạt động tốt hơn. Difficulty
label có tương quan với question type: nhóm hard chứa nhiều mẫu
counterexample, chain hoặc cấu trúc có gold answer dễ khớp theo template,
trong khi nhóm medium chứa phần lớn reverse reasoning. Vì vậy, kết quả theo
difficulty chịu ảnh hưởng của thành phần dữ liệu và cần được phân tích cùng
question type.

\subsection{Kết quả trên counterfactual subset}

Trên \(20\) mẫu có
\texttt{requires\_counterfactual=true}, hệ thống đạt:

\begin{itemize}
    \item Verification Accuracy: \(100{,}00\%\);
    \item Answer Token F1: \(89{,}16\%\);
    \item Citation F1: \(67{,}50\%\);
    \item Rule Recall@5: \(82{,}50\%\);
    \item Event Recall@5: \(88{,}33\%\);
    \item Top-1 Exact Path: \(55{,}00\%\);
    \item Oracle Exact Path: \(95{,}00\%\).
\end{itemize}

Khoảng cách giữa Oracle và Top-1 Exact Path trên subset này là \(40\) điểm
phần trăm. Mặc dù verifier phân loại đúng toàn bộ direct-edge negative
claims, path ranking vẫn chưa luôn chọn gold path.

Như đã nêu tại Phần~\ref{sec:data}, các mẫu này chủ yếu kiểm tra việc bác bỏ
cạnh trực tiếp. Do đó, kết quả \(100\%\) không được sử dụng để kết luận rằng
LegalSCM đã giải quyết hoàn chỉnh hard-intervention reasoning.

\subsection{Paired ablation}

Bảng~\ref{tab:paired-ablation} so sánh structural SCM và path ablation trong
paired setup.

\begin{table*}[t]
    \centering
    \caption{Paired comparison giữa structural SCM và path ablation.}
    \label{tab:paired-ablation}
    \begin{tabular}{lrrr}
        \hline
        \textbf{Metric} &
        \textbf{Structural SCM} &
        \textbf{Path ablation} &
        \textbf{Chênh lệch} \\
        \hline
        Rule Recall@5 &
        \(69{,}93\%\) & \(69{,}93\%\) & \(0{,}00\) \\
        Event Recall@5 &
        \(74{,}66\%\) & \(74{,}66\%\) & \(0{,}00\) \\
        Top-1 Exact Path &
        \(35{,}20\%\) & \(35{,}20\%\) & \(0{,}00\) \\
        Oracle Exact Path &
        \(71{,}20\%\) & \(71{,}20\%\) & \(0{,}00\) \\
        Verification Accuracy &
        \(92{,}00\%\) & \(92{,}00\%\) & \(0{,}00\) \\
        Verification Macro-F1 &
        \(97{,}73\%\) & \(97{,}73\%\) & \(0{,}00\) \\
        Answer Token F1 &
        \(68{,}69\%\) & \(68{,}69\%\) & \(0{,}00\) \\
        Answer ROUGE-L F1 &
        \(67{,}67\%\) & \(67{,}67\%\) & \(0{,}00\) \\
        Citation F1 &
        \(62{,}86\%\) & \(62{,}86\%\) & \(0{,}00\) \\
        Counterfactual Accuracy &
        \(100{,}00\%\) & \(100{,}00\%\) & \(0{,}00\) \\
        \hline
        Thời gian trung bình/mẫu &
        \(3{,}890\) s & \(1{,}205\) s & \(+2{,}685\) s \\
        Thời gian trung vị/mẫu &
        \(2{,}269\) s & \(1{,}220\) s & \(+1{,}049\) s \\
        Batch wall-clock time &
        \(1\,284{,}455\) s & \(450{,}429\) s & \(+834{,}026\) s \\
        \hline
    \end{tabular}
\end{table*}

Hai verification engine tạo ra các quality metric giống hệt nhau trên
benchmark hiện tại. Vì retrieval được chạy trước Step 4, việc hai cấu hình
có cùng retrieval và path metric là kết quả được kỳ vọng. Quan trọng hơn,
verification, answer và citation metric cũng không thay đổi.

Về runtime, tỷ lệ thời gian trung bình trên mỗi mẫu là:

\begin{equation}
\frac{
    3{,}890
}{
    1{,}205
}
\approx
3{,}23.
\end{equation}

Như vậy, structural SCM chậm hơn khoảng \(3{,}23\) lần theo per-sample
runtime. Theo batch wall-clock time, tỷ lệ là:

\begin{equation}
\frac{
    1\,284{,}455
}{
    450{,}429
}
\approx
2{,}85.
\end{equation}

Chênh lệch giữa hai tỷ lệ có thể đến từ chi phí khởi tạo, cache, checkpoint
và ghi artifact không được phân bổ giống nhau trong per-sample runtime.

Kết quả paired được xem là một \emph{null quality ablation}: structural SCM
không làm tăng quality metric trên benchmark hiện tại. Lợi ích của
structural SCM nằm ở semantics hard intervention và trace chi tiết hơn,
nhưng các lợi ích này chưa được benchmark hiện tại đánh giá trực tiếp.

\subsection{Trả lời các câu hỏi nghiên cứu}

\paragraph{RQ1 --- Chất lượng retrieval.}

Hệ thống đạt Rule Recall@5 \(69{,}93\%\) và Event Recall@5 \(74{,}66\%\).
Retrieval có coverage tương đối cao nhưng precision thấp do giữ nhiều
candidate rule.

\paragraph{RQ2 --- Chất lượng causal path.}

Top-1 Exact Path chỉ đạt \(35{,}20\%\), trong khi Oracle Exact Path đạt
\(71{,}20\%\). Kết quả xác định path ranking là điểm nghẽn lớn hơn so với
khả năng tạo candidate path trong nhiều trường hợp.

\paragraph{RQ3 --- Chất lượng xác minh.}

Query-aware verifier đạt Accuracy \(92{,}00\%\), phân loại đúng toàn bộ
\(20\) direct-edge negative claims và chuyển \(20\) gold supported sample
thành uncertain.

\paragraph{RQ4 --- Chất lượng câu trả lời.}

Answer Token F1 đạt \(68{,}69\%\), ROUGE-L F1 đạt \(67{,}67\%\) và Citation
F1 đạt \(62{,}86\%\). Reverse reasoning là nhóm có hiệu năng thấp nhất.

\paragraph{RQ5 --- Ảnh hưởng của verification engine.}

Structural SCM và path ablation có quality metric giống nhau, nhưng
structural SCM chậm hơn khoảng \(3{,}23\) lần trên mỗi mẫu. Benchmark hiện
tại chưa đủ nhạy để đo lợi ích accuracy của hard-intervention semantics.
